\documentclass[10pt,a4paper]{amsart}
\usepackage[margin=19mm]{geometry}
\usepackage{amsmath,amssymb,mathtools}
\usepackage[scr=rsfs,cal=euler]{mathalfa}
\usepackage{xfrac}
\usepackage[unicode,hidelinks,bookmarks=false]{hyperref}
\newcommand{\ie}{i.e.\ }
\newcommand{\cf}{cf.\ }
\newcommand{\resp}{respectively\ }
\newcommand{\Subsection}{\subsection}
\providecommand{\nobreakdash}{\nobreak\hbox{-}}
\setlength{\emergencystretch}{2em}
\multlinegap0pt
\usepackage{amssymb}
\usepackage{mathtools}
\newtheorem{proposition}{Proposition}
\newcommand{\LM}{\mathscr{LM}}
\newcommand{\calM}{\mathcal{M}}
\title{Two-place Laplacian matching root integral variations are impossible}
\author{Sebastian M. Cioab\u{a}, Lele Liu, Yi Wang}
\date{Source: arXiv:2605.01760v1 (CC BY 4.0)}
\begin{document}
\maketitle

\noindent\emph{Source and licence.} Two-place Laplacian matching root integral variations are impossible. Version arXiv:2605.01760v1 (CC BY 4.0), distributed by arXiv under the Creative Commons Attribution 4.0 International licence, http://creativecommons.org/licenses/by/4.0/, as recorded in the arXiv licence metadata for this exact version and shown on the abstract page https://arxiv.org/abs/2605.01760v1. Full licence notice: \texttt{licenses/arxiv-2605.01760-CC-BY-4.0.txt}.\par\smallskip

\noindent\textit{Editorial scope.} Counted unit: the article's proposition prop:powersums (source0.tex lines 145--154). The article's complete original proof of it is delivered with it (source0.tex lines 156--246, one \texttt{proof} environment of 91 lines). No mathematics is written, repaired or re-derived: the statement and the proof are the article's own lines, in the article's own notation, and no retained line is substituted or re-flowed.  Retained uncounted as the source-local context the proof needs: lines 86--88.  Nothing was omitted from any retained span, so the package carries no omission record.  No retained line contains a citation, so the wrapper carries no bibliography and no bibliography entry.  No retained line contains a figure environment, an image-inclusion command or drawing code, so no image is delivered and the package declares no asset dependency.  Editorial formatting only, in addition: an AMS \texttt{amsart} preamble that restates the article's own declarations and macros used by the retained lines, namely the environments \texttt{proposition} on separate counters, together with the standard packages those lines need.  The retained environments print in this document as Proposition 1 in order of appearance, whereas the article prints the same items with the numbering of its own full document.  The complete native source of the article as distributed in the arXiv e-print for this exact version is bundled unchanged as \texttt{sources/199657/source0.tex}.
\begin{equation}\label{eq:LMdef}
\LM_G(x)=\sum_{M\in \calM(G)}(-1)^{|M|}\prod_{v\in V(G)\setminus V(M)}(x-d_G(v)).
\end{equation}

\begin{proposition}\label{prop:powersums}
Let $G$ be a graph. Then
\[
p_4(G)=A_4+4A_3+2A_2+4B-2|E(G)|,
\]
and
\[
p_5(G)=A_5+5A_4+5A_3-5A_2+5C+10B.
\]
\end{proposition}

\begin{proof}
Write $\LM_G(x)=x^n-b_1x^{n-1}+b_2x^{n-2}-b_3x^{n-3}+b_4x^{n-4}-b_5x^{n-5}+\cdots +(-1)^nb_n$.

Denote by $E_r$ the $r$-th elementary symmetric polynomial in the degree sequence \\
$\{d(v):v\in V(G)\}$. By expanding \eqref{eq:LMdef}, one sees that only matchings of 
size at most $2$ contribute to $b_j$ for $1\le j\le 5$. A direct coefficient computation yields
\begin{align*}
b_1 &= E_1,\\
b_2 &= E_2-|E(G)|,\\ 
b_3 &= E_3-\sum_{xy\in E(G)} \bigl(E_1-d(x)-d(y)\bigr),\\
b_4 & = E_4-\sum_{xy\in E(G)}\bigl(E_2-E_1(d(x)+d(y))
+d(x)^2+d(y)^2+d(x)d(y)\bigr) + \phi_2(G), \\
b_5 & = E_5-\sum_{xy\in E(G)} \!\bigl(E_3-E_2(d(x)+d(y))
+E_1\bigl(d(x)^2+d(x)d(y)+d(y)^2\bigr)\bigr) \\
& \quad+\sum_{xy\in E(G)}\! \bigl(d(x)^3+d(x)^2d(y)+d(x)d(y)^2+d(y)^3\bigr) 
+\sum_{M\in \calM_2(G)}\!\! \bigg(E_1-\sum_{z\in V(M)} d(z)\bigg),
\end{align*}
where $\calM_2(G)$ is the set of $2$-matchings of $G$.

Now, by the definition of $A_r$, we have
\[
A_2 = \sum_{xy\in E(G)}\! \bigl(d(x)+d(y)\bigr),~~~ 
A_3 = \sum_{xy\in E(G)}\! \bigl(d(x)^2+d(y)^2\bigr),~~~
A_4 = \sum_{xy\in E(G)}\! \bigl(d(x)^3+d(y)^3\bigr), 
\]
and
\begin{align*}
\sum_{M\in \calM_2(G)} \sum_{z\in V(M)} d(z)
& = \sum_{xy\in E(G)} \bigl(d(x) + d(y)\bigr) \bigl(|E(G)|-d(x)-d(y)+1\bigr) \\
& = (|E(G)|+1) A_2 - A_3 - 2B.
\end{align*}
Set $m:=|E(G)|$. Using the identities above, together with
\[
\phi_2(G)=\binom{m}{2}-\sum_{v\in V(G)}\binom{d(v)}{2}
=\frac{m^2+m-A_2}{2},
\]
the coefficient formulas become 
\begin{align*}
b_1 &= E_1,\\ 
b_2 &=E_2-m,\\ 
b_3 &=E_3-mE_1+A_2,\\
b_4 & = E_4-mE_2+E_1A_2-A_3-B+\phi_2(G), \\
b_5 & = E_5-mE_3+E_2A_2-E_1(A_3+B)+A_4+C \\
& \quad +E_1\phi_2(G)-(m+1)A_2+A_3+2B.
\end{align*}

Now let $\lambda_1,\dots,\lambda_n$ be the roots of $\LM_G(x)$. 
Newton's identities imply that
\begin{align*}
p_1 &= b_1,\\ 
p_2 &=b_1^2-2b_2,\\ 
p_3 &=b_1^3-3b_1b_2+3b_3,\\
p_4 & = b_1^4-4b_1^2b_2+2b_2^2+4b_1b_3-4b_4, \\
p_5 & = b_1^5-5b_1^3b_2+5b_1b_2^2+5b_1^2b_3-5b_2b_3-5b_1b_4+5b_5.
\end{align*}

Substituting the expressions for $b_1,b_2,b_3,b_4$ into the formula for $p_4$, we get that
\begin{align*}
p_4 & = E_1^4-4E_1^2(E_2-m)+2(E_2-m)^2+4E_1(E_3-mE_1+A_2) \\
& \quad -4\bigl(E_4-mE_2+E_1A_2-A_3-B+\phi_2(G)\bigr) \\
& = E_1^4-4E_1^2E_2+2E_2^2+4E_1E_3-4E_4+4A_3+4B+2m^2-4\phi_2(G).
\end{align*}
We now apply Newton's identities to the degree sequence
$d(v_1),\dots,d(v_n)$, whose elementary symmetric polynomials are precisely
$E_1,E_2,\dots, E_n$, and we get that 
\begin{equation*}
A_4=E_1^4-4E_1^2E_2+2E_2^2+4E_1E_3-4E_4
\end{equation*}
and 
\[
p_4 = A_4+4A_3+4B+2m^2-4\phi_2(G) 
= A_4+4A_3+2A_2+4B-2m.
\]

Similarly, substituting the expressions for $b_1,\dots,b_5$ into the formula for $p_5$ gives
\begin{align*}
p_5 & = E_1^5-5E_1^3(E_2-m)+5E_1(E_2-m)^2+5E_1^2(E_3-mE_1+A_2) \\
& \quad -5(E_2-m)(E_3-mE_1+A_2) 
-5E_1\bigl(E_4-mE_2+E_1A_2-A_3-B+\phi_2(G)\bigr) \\
& \quad +5\bigl(E_5-mE_3+E_2A_2-E_1(A_3+B)+A_4+C \\
& \quad +E_1\phi_2(G)-(m+1)A_2+A_3+2B\bigr) \\
& = E_1^5-5E_1^3E_2+5E_1E_2^2+5E_1^2E_3-5E_2E_3-5E_1E_4+5E_5 \\
& \quad +5A_4+5A_3-5A_2+10B+5C.
\end{align*}
Applying Newton's identities to the degree sequence again, we have
\[
A_5=E_1^5-5E_1^3E_2+5E_1E_2^2+5E_1^2E_3-5E_2E_3-5E_1E_4+5E_5.
\]
Therefore, $p_5=A_5+5A_4+5A_3-5A_2+10B+5C$.
This completes the proof.
\end{proof}

\end{document}
